\Subsection*{4) Proof of 2)\kern2pt b)}

4) 1. Case where $\Phi$ is a gerbe.

One can find a surjective \'etale morphism of finite type $p
\colon X_1 \to X$ such that $p^*\Phi$ is a trivial gerbe. By
descent, as in 3)\kern2pt2, one sees that it suffices to prove the
theorem for $X_1$, $X_1\times_X X_1$ and
% original p. 403
$X_1\times_X X_1\times_X X_1$. One is therefore reduced to the case where
$\Phi$ is the gerbe of torsors under a constructible sheaf of groups
$F$ whose fibers are of order prime to the
residue characteristics of~$S$.

By SGA~4 IX~2.14, one can find a finite family of
finite morphisms $p_i \colon Z_i \to X$ and, for each $i$, a
finite constant sheaf of groups $C_i$, of order prime to the
residue characteristics of~$S$, such that one has a
monomorphism
$$
j\colon F \to \prod_{i} p_{i*}C_i=G \quoi.
$$
Let $\Phi_i$ be the stack of torsors under $C_i$ and $\Psi$ the stack of
torsors under $G$. It follows from 2)\kern2pt a) that, at the cost of
restricting $S$ to a nonempty open subset, one may assume that the
$(C_i,fp_i)$ are cohomologically proper relative to~$S$ in
dimension $\le 1$, and that the stacks $f_*p_{i*}\Phi_i$ are
$1$-constructible. It then follows from~\Ref{XIII.1.9} that the
$(p_{i*}C_i,f)$ are cohomologically proper relative to~$S$
in dimension $\le 1$; hence the same holds for~$(G,f)$. Moreover,
since $p_{i*}\Phi_i$ is equivalent to the stack of torsors under the
group $p_{i*}C_i$ (SGA~4 VIII~5.8), one sees that $f_*\Psi$ is
$1$-constructible.

Since $\R^1f_*G$ is constructible, one can find a sheaf
representable by an \'etale $Y$-scheme of finite type $T$ and
an epimorphism
$$
a \colon T \to \R^1f_*G
$$
(SGA~4 IX~2.7); moreover one may assume that the image of the identity
section of~$T(T)$ is defined by a torsor $Q$ on~$X\times_Y
T=X_T$, with group $G|X_T$. Let $f_T \colon X_T \to T$ be the canonical
morphism, $F_T=F|X_T$, etc. By 1)\kern2pt b) one may assume,
at the cost of restricting $S$ to a nonempty open subset, that
$(Q/F_T,f_T)$ is cohomologically proper relative to~$S$ in
dimension $\le 0$ and that $f_{T*}(Q/F_T)$ is constructible. It
then follows from~\Ref{XIII.3.1.2} that $f_*\Psi$ is constructible.

Let us show that $(F,f)$ is cohomologically proper relative to
% original p. 404
$S$ in dimension $\le 1$. By \Ref{XIII.1.13}~2) it suffices to
prove that, for every scheme $Y_1$ \'etale over~$Y$ and every
torsor $Q_1$ on~$X_1=X\times_Y Y_1$, if $f_1 \colon X_1 \to Y_1$ is
the canonical morphism, then $(Q_1/F_1,f_1)$ is cohomologically
proper relative to~$S$ in dimension $\le 0$. Now, by
definition of~$T$, $Q_1$ is, locally for the \'etale topology
of~$Y_1$, an inverse image of~$Q$, which proves our
reduction.

4) 2. General case.\par
Using lemma~\Ref{XIII.6.1.1}, 4)\kern2pt1 and 1)\kern2pt a), one sees that,
at the cost of restricting $S$ to a nonempty open subset, one may
assume $S(f_*\Phi)$ constructible and $(S\Phi,f)$ cohomologically
proper relative to~$S$ in dimension $\le 0$. One can then
find a sheaf representable by an \'etale $Y$-scheme
of finite type $T$ and an epimorphism
$$
a \colon T \to S(f_*\Phi)\quoi.
$$
One takes up again the notation of 4)\kern2pt1, and one sets moreover $Z=T\times_Y T$,
$X_Z=X\times_Y Z$, and one denotes by $f_Z \colon X_Z \to Z$ the canonical
morphism. One may assume $T$ chosen so that the image $q$ of the
identity section of~$T(T)$ under $a$ is defined by an object~$p$
of~$(f_*\Phi)_T=\Phi_{X_T}$. Let $p_1$ and $p_2$ (\resp~$q_1$ and
$q_2$) be the inverse images of~$p$ (\resp~$q$) under the two projections
of~$Z$ into $T$. One may assume, at the cost of restricting $S$ to
a nonempty open subset, that $(\SheafAut_{X_T}(p),f_T)$ is
cohomologically proper relative to~$S$ in dimension $\le 1$,
that $(\SheafHom_{X_Z}(p_1,p_2),f_Z)$ is cohomologically proper
relative to~$S$ in dimension $\le 0$, and that
$f_{T*}(\SheafAut_{X_T}(p))$, $f_{Z*}(\SheafHom_{X_Z}(p_1,p_2))$ are
constructible ( a)\kern2pt1) and 4)\kern2pt1)).

From this one first deduces that, for every scheme $Y_1$ \'etale over~$Y$ and every pair of objects $y$, $y_1$ of~$(f_*\Phi)_{Y_1}$, the
sheaf $\SheafAut_{Y_1}(y)$ (\resp $\SheafHom_{Y_1}(y,y_1)$) is
constructible, such a sheaf being, locally for the \'etale
topology of~$Y_1$, an inverse image of~$f_{T*}(\SheafAut_{X_T}(p))$
(\resp of~$f_{Z*}(\SheafHom_{X_Z}(p_1,p_2))$).

It remains to prove that $(\Phi,f)$ is cohomologically proper
relative
% original p. 405
to~$S$ in dimension $\le 1$. For this it suffices to show that,
for every $S'$-scheme $S$ and every geometric point
$\overline{y}'$ of~$Y$, if one denotes by $\overline{Y}$ the strict
localization of~$Y$ at $y'$, $\overline{X}=X \times_Y \overline{Y}$, etc.,
then the canonical functor
$$
\overline{\varphi} \colon \Phi(\overline{X}) \to
\Phi'(\overline{X'})
$$
is an equivalence of categories.

Let us show that $\overline{\varphi}$ is fully faithful. Let
$x$, $y \in \Phi(\overline{X})$, let $x'$, $y'$ be their images in
$\Phi'(\overline{X'})$, and let us show that the canonical morphism
\begin{equation*}
\label{eq:XIII.4.*} \tag{$*$} {\Hom_{\overline{X}}(x,y)
\to \Hom_{\overline{X'}}(x',y') }
\end{equation*}
is bijective. By definition of~$T$ there exist two morphisms from
$\overline{Y}$ to $T$ such that $x$ and~$y$ are the inverse images
of~$p$ under these two morphisms. This amounts to saying that there is a
morphism $\overline {Y} \to Z$, whence a morphism $h\colon
\overline{X} \to X_Z$ such that one has
$$
h^*(p_1)=x \qquad h^*(p_2)=y.
$$
Consequently one has a canonical isomorphism
$$
\SheafHom_{\overline{X}}(x,y)=h^*(\SheafHom_{X_Z}(p_1,p_2) )\quoi.
$$
but, taking into account the fact that $(\SheafHom_{X_Z}(p_1,p_2),f_Z)$ is
cohomologically proper relative to~$S$ in dimension $\le 0$,
one sees that the same holds for
$(\SheafHom_{\overline{X}}(x,y),\overline{f})$, which proves that the
morphism \eqref{eq:XIII.4.*} is bijective.

Let us show that $\overline{\varphi}$ is essentially surjective. Let
$x'\in \Phi'(\overline{X'})$. Since $(S\Phi,f)$ is cohomologically
proper relative to~$S$ in dimension $\le 0$, the canonical
morphism
$$
\H^0(\overline{X},\overline{S\Phi}) \to
\H^0(\overline{X'},\overline{S\Phi}')
$$
is bijective. Let $G'$ be the maximal subgerbe of~$\Phi'$
generated by $x'$; there then exists a maximal subgerbe $G$ of
$\overline{\Phi}$ whose inverse image on~$\overline{X'}$ is
$G'$. By 4)\kern2pt1, $(G,\overline{f})$ is cohomologically proper
relative to~$S$ in dimension $\le 1$; consequently the canonical
functor
$$
G(\overline{X}) \to G'(\overline{X'})
$$
% original p. 406
is an equivalence of categories, which proves the existence
of an element $x$ of~$\Phi(\overline{X})$ whose image in
$\Phi'(X')$ is isomorphic to $x'$, and completes the
proof of the theorem.

\begin{sublemma}
\label{XIII.3.1.1}
Let $S$ be a locally noetherian scheme and
$$
\xymatrix@C=.5cm{ \Phi \ar[r]^p & \Phi_1 \ar@<2pt>[rr]^{p_1,p_2}
\ar@<-2pt>[rr] && \Phi_2 }
$$
an exact diagram of stacks on~$S$ \eqref{XIII.1.10.1}. If $\Phi_1$
is constructible, so is $\Phi$. If, for every
scheme $S'$ \'etale over~$S$ and every pair
of objects
$x_1,~y_1$ of~$(\Phi_1)_{S'}$, the sheaf $\SheafHom_{S'}(x_1,y_1)$
is constructible, then, for every pair of objects $x,~y$ of
$\Phi_{S'}$, the same holds for~$\SheafHom_{S'}(x,y)$. Suppose
that $\Phi_1$ is $1$-constructible (\Ref{XIII.0}) and that $\Phi_2$ is
constructible; then $\Phi$ is $1$-constructible.
\end{sublemma}

For every scheme $S'$ \'etale over~$S$ and every object $x$ of
$\Phi_{S'}$, one has a monomorphism
$$
\SheafAut_{S'}(x)\to\SheafAut_{S'}(p(x)).
$$
It therefore follows from SGA~4 IX~2.9 that, if $\Phi_1$ is
constructible, so is $\Phi$, and the second
assertion of the lemma is proved in the same way.

Suppose now $\Phi_1$ $1$-constructible and $\Phi_2$
constructible. The morphism $p$ induces on the sheaves of
maximal subgerbes a morphism
$$
\varphi\colon S\Phi\to S\Phi_1.
$$
Let $G$ be the image of~$S\Phi$ under $\varphi$; by SGA~4 IX~2.9,
$G$ is a constructible sheaf. One can therefore find a sheaf
representable by an \'etale $S$-scheme of finite type $T$ and
an epimorphism
$$
a\colon T\to G
$$
(SGA~4 IX~2.7).
% original p. 407
Moreover one can choose $T$ so that the image $y$ of the identity
section of~$T(T)$ under $a$ is defined by an object $x_1$ of
$(\Phi_1)_T$ of the form $x_1=p(x)$, where $x\in\Phi_T$.

It suffices to show that, for every point $s$ of~$S$, there exists a
nonempty open subset~$U$ of~$\overline{\{ s\} }$ such that $S\Phi|U$ is
locally constant constructible. Let $s\in S$, let $\overline{s}$ be a
geometric point over~$s$ and $\overline{y}_1,\dots,\overline{y}_n$ the elements of~$G_{\overline{s}}$. By
definition of~$T$ there exist morphisms
$h_i\colon\overline{s}\to T$ such that one has
$h_i^*(y)=\overline{y}_i$. Let $S'$ be the fiber product over~$S$ of
$n$ schemes isomorphic to $T$, $h\colon\overline{s}\to S'$ the
fiber product of the $h_i$, $y_i$ (\resp $x_i$) the inverse image of
$y$ (\resp $x$) under the $i$th projection of~$S'$ onto~$T$. Let
$F_i$ be the subsheaf of~$S\Phi |S'$ that is the inverse image of~$y_i$,
and let us show that the $F_i$ are constructible. The sheaf $F_i$ is
a quotient of the sheaf $F'_i$ such that, for every scheme $S''$
\'etale over~$S'$, one has
\begin{multline*}
F'_i(S'') = \{\text{classes mod.\ isomorphism of objects $z$ of
$\Phi_{S''}$ equipped}\\
\text{ with an isomorphism $i\colon p(z)\simeq p(x_i|S'')$}\}.
\end{multline*}
It suffices to show that the $F_i'$ are constructible. Now, if one
sets $z_i=p_1p(x_i)$, $z_i'=p_2p(x_i)$, one has a monomorphism
$$
\Psi_i\colon F'_i\to\SheafIsom_{S'}(z_i,z_i'),
$$
obtained by associating, with every scheme $S''$ \'etale over~$S'$ and
every object $z$ of~$\Phi_{S''}$ such that one has an isomorphism
$i\colon p(z)\simeq p(x_i|S'')$, the isomorphism from~$z_i$ to $z_i'$
defined by the condition that the diagram
$$
\xymatrix@C=1.2cm{
p_1p(z) \ar[r]^-{p_1(i)}\ar[d]_-j& z_i\ar[d]\\ p_2p(z) \ar[r]^-{p_2(i)}& z'_i
}
$$
be commutative ($j$ is the canonical morphism associated with the
exact diagram).
% original p. 408

The morphism $\Psi_i$ is injective because to say that two objects $z,~z'$
such that one has isomorphisms $i\colon p(z)\isomto
p(x_i|S''),~i'\colon p(z')\isomto p(x_i|S'')$ define the
same element of~$\SheafIsom_{S''}(z_i,z'_i)$ amounts to
saying that one has $p_1(i'^{-1}i)=p_2(i'^{-1}i)$, \ie that $i'^{-1}i$
comes from an isomorphism $z\to z'$. The sheaf $F'_i$, being a
subsheaf of~$\SheafIsom_{S'}(z_i,z'_i)$, is constructible.

Since one can find a nonempty open subset~$U$ of~$\overline{\{ s\}}$
such that $y_1|U,\dots,y_n|U$ generate $G$, it follows from
lemma~\Ref{XIII.6.1.2} below that $S\Phi |U$ is constructible,
hence, at the cost of shrinking $U$, $S\Phi |U$ is locally constant
constructible.

\begin{sublemma}
\label{XIII.3.1.2}
Let $S$ be a locally noetherian scheme, $f\colon X\to S$
a morphism, $F\to G$ a monomorphism of sheaves of groups on~$X$, $\Psi$ (\resp $\Psi_1$) the stack of torsors under $F$
(\resp under $G$), $\Phi=f_*\Psi$, $\Phi_1=f_*\Psi_1$. Suppose
given a sheaf on~$S$ representable by an \'etale $S$-scheme
of finite type $T$ and a surjective morphism
$$
a\colon T\to S\Phi_1\simeq\R^1f_*G,
$$
such that there exists a torsor $Q$ on~$X_T=X\times_ST$ with group
$G|X_T$ which defines in $\R^1f_*G(T)$ the image under $a$ of the
identity section of~$T(T)$. Let $f_T\colon X_T\to T$ be the canonical
morphism and set $F_T=F|X_T$. Suppose that $\Phi_1$ is
$1$-constructible and that $f_{T*}(Q/F_T)$ is constructible. Then
$\Phi$ is $1$-constructible.
\end{sublemma}

It suffices indeed to copy the proof
of~\Ref{XIII.3.1.1}, the fact that one has morphisms $\Psi_i$
being replaced by the fact that one has isomorphisms
$$
F'_i\isomto f_*(Q/F_T)|S'.
$$

\begin{subremark}
\label{XIII.3.1.3}
Suppose
% original p. 409
that $k=\kres(s)$ is of characteristic zero; then the
schemes of finite type over~$\overline{k}$ of dimension $\leq\dim X$
are strongly desingularizable, and the proof
of~\Ref{XIII.3.1} makes it possible to prove the following results:
\begin{enumerate}
\item[a)] There exists a nonempty open subset~$S_1$ of~$S$ such that, for every
scheme $S'$ over~$S_1$ whose maximal points are of
characteristic zero and every locally
constant constructible sheaf of sets $F$ on~$X'=X\times_SS'$, $(F,f')$ is
cohomologically proper relative to~$S'$ in dimension $\leq
0$.

\item[b)] If all the residue characteristics of~$S$ are
zero, there exists a nonempty open subset~$S_1$ of~$S$ such that, for every
scheme $S'$ over~$S_1$ and every locally constant
constructible sheaf of groups $F$ on~$X'$, $(F,f')$ is
cohomologically proper relative to~$S'$ in dimension $\leq1$.
\end{enumerate}

It suffices indeed to copy the proof of
\Ref{XIII.3.1}.2)\kern2pt a). Proposition~\Ref{XIII.2.7} used in
3)\kern2pt 4 applies to the case of a locally constant sheaf $F$, because,
taking up again the notation of 3)\kern2pt4, every torsor under $F$ is
tamely ramified on~$P$ relative to~$S$, since
all the residue characteristics of~$S$ are zero.
\end{subremark}

\begin{corollary}
\label{XIII.3.2}
Let $k$ be a field of characteristic $p\geq 0$, $p'$ the set
of prime numbers distinct from~$p$, $f\colon X\to k$ a coherent
morphism.
\begin{enumerate}
\item[1)] For every sheaf of sets $F$, $(F,f)$ is
cohomologically proper in dimension $\leq 0$.
\item[2)] Suppose that one of the following two conditions is satisfied:
\begin{enumerate}
\item[a)] $f$ is of finite type and the schemes of finite type of
dimension $\leq\dim X$ over an algebraic closure of~$k$ are
strongly de\-sin\-gu\-lar\-iz\-able.
\item[b)] The schemes of finite type over an algebraic
closure of~$k$ are
% original p. 410
strongly desingularizable.
\end{enumerate}
Then, for every sheaf of ind-$p'$-groups $F$, $(F,f)$ is
cohomologically proper in dimension $\leq 1$.
\end{enumerate}
\end{corollary}

Let $F$ be a sheaf of sets (\resp of
ind-$p'$-groups). By SGA~4 IX~2.7.2 one can write
$F$ as a filtered inductive limit
$$
F=\varinjlim F_i,
$$
where the $F_i$ are constructible sheaves of sets (\resp of
ind-$p'$-groups). Since $f$ is coherent, $f_*$
(\resp $\R^1f_*$) commutes with inductive limits (SGA~4 VII~3.3). If
one knows that the $(F_i,f)$ are cohomologically proper in
dimension $\leq 0$ (\resp that every sheaf of sets is
cohomologically proper in dimension $\leq 0$ and that the $(F_i,f)$
are cohomologically proper in dimension $\leq 1$), the same will
hold for~$(F,f)$. One may therefore assume $F$ constructible.

If one assumes $f$ of finite type (\resp satisfying a)), the
proposition follows from \Ref{XIII.3.1}~1)\kern2pt b) (\resp from
\Ref{XIII.3.1}~2)\kern2pt b)). Let us now prove~\Ref{XIII.3.2} when one no longer
assumes $f$ of finite type. For every scheme $S'$ over
$k$ and every geometric point $\overline{s}$ of~$S'$, one
denotes by $\overline{k}$ (\resp $\overline{S'}$) the strict localization of
$k$ at $\overline{s}$ (\resp of~$S'$ at $\overline{s}$),
by $\overline{X}$ the inverse image of~$X$ over~$\overline{k}$, and one
considers the cartesian square
$$
\xymatrix@C=1.2cm{
\overline{X}\ar[d]_-{\overline{f}}&\ar[l]_g \overline{X'}\ar[d]^-{\overline{f}'}\\ \overline{k}&\ar[l]\overline{S'}
}
$$
It suffices to prove that, for every $S'$ and every $\overline{s}$, the
canonical morphism
$$
\H^0(\overline{X},\overline{F})\to
\H^0(\overline{X'},\overline{F'})\qquad
\text{(\resp $\H^1(\overline{X},\overline{F})\to
\H^1(\overline{X'},\overline{F'})$)}
$$
is an isomorphism. It suffices to show that one has the relations
% original p. 411
\begin{equation*}
\label{eq:XIII.3.*}
\tag{$*$} \overline{F}\simeq g_*g^*\overline{F}\qquad
\text{(\resp $\R^1g_*(g^*\overline{F})=0$).}
\end{equation*}
Now, in this form, the question is local on~$X$ for the \'etale
topology. One may therefore assume $X$ affine; by passage to the
limit one may assume $X$ of finite type over~$k$. One then knows that
$(F,f)$ is cohomologically proper in dimension $\leq 0$
(\resp $\leq 1$) and that the same holds when $X$ is replaced by
a scheme \'etale of finite type over~$X$, which proves
\eqref{eq:XIII.3.*}.

\begin{theorem}
\label{XIII.3.3}
Let $S$ be an irreducible scheme with generic point
$s$, $f\colon X\to S$ a morphism of finite
presentation. Suppose that the schemes of finite type of dimension
$\leqslant\dim X_s$ over an algebraic closure $\bar{k}$ of
$k$ are desingularizable \textup{(EGA~IV 7.9.1)}. Then, if $\LL$
denotes the set of prime numbers distinct from the
residue characteristics of~$S$, one can find a nonempty open
subset~$S_1$ of~$S$ such that the morphism $f|S_1$ is
universally locally $1$-aspherical for~$\LL$.
\end{theorem}

One may assume $S$ integral and $X$ reduced (SGA~4
VIII~1.1). By passage to the limit one may assume $S$
noetherian. Moreover, to prove the theorem, it
suffices to do so after a finite extension $S_1\to S$, where $S_1$
is an integral scheme and where $S_1\to S$ is a composite
of \'etale extensions and of surjective radicial extensions.

Let us first show that, at the cost of restricting $S$ to a nonempty open
subset, $f$ is universally locally $0$-acyclic. At the cost of
making a radicial extension of~$\kres(s)$, one may assume that the
morphism $(X_s)_\red\to s$ is separable; one may therefore assume,
at the cost of restricting $S$ to a nonempty open subset and of making
a surjective radicial extension of~$S$, that the morphism $f$ is
flat, with separable geometric fibers (EGA~IV
12.1.1), which implies that $f$ is universally
$0$-acyclic (SGA~4 XV~4.1).

Let us show that, at the cost of restricting $S$ to a nonempty open subset,
$f$ is universally locally $1$-aspherical for
$\LL$. At the cost of making a
% original p. 412
finite extension of~$\kres(s)$, which is permissible because one can
consider it as a composite of an \'etale extension and of a
radicial extension, one can find a surjective proper morphism
$p_s\colon Y_s\to X_s$, where $Y_s$ is a scheme smooth over~$S$,
of the same dimension as $X_s$, and one may assume that $p_s$
comes from a surjective proper morphism $p\colon Y\to X$, where $Y$
is a scheme smooth over~$S$ (EGA~IV 9.6.1 and 12.1.6). It suffices to
show that, at the cost of restricting $S$ to a nonempty open subset,
for every diagram with cartesian squares
$$
\xymatrix{ S''\ar[d]_i & X ''\ar[d]^j\ar[l]_{f''} \\ S'\ar[d] &
X'\ar[d]\ar[l]_{f'} \\ S & X, \ar[l]_f }
$$
where $i$ is \'etale of finite presentation, and for every
sheaf of ind-$\LL$-groups $F$ on~$S''$, if $\Phi$ is the stack
of torsors under $F$, then the canonical morphism
$$
{f'}^*i_*\Phi \to j_*{f''}^*\Phi
$$
is an equivalence. Let $Z=Y\times_{X}Y,
T=Y\times_{X}Y\times_{X}Y$. One has in a natural way a commutative
diagram
$$
\xymatrix{ S'' \ar[d]_i & X''\ar[l]_{f''}\ar[d]_j &
Y''\ar[l]_{p''}\ar[d] & Z''\ar[d]\ar@<-.4ex>[l]\ar@<.4ex>[l] &
T''\ar[d]\ar[l]\ar@<-.8ex>[l]\ar@<.8ex>[l] \\ S' \ar[d] & X'
\ar[l]_{f'}\ar[d] & Y'\ar[d]\ar[l]_{p'}&
Z'\ar[d]\ar@<-.4ex>[l]\ar@<.4ex>[l] &
T'\ar[d]\ar[l]\ar@<-.8ex>[l]\ar@<.8ex>[l] \\ S & X\ar[l]_f & Y\ar[l]_p
& Z'\ar@<-.4ex>[l]\ar@<.4ex>[l] & T'\ar[l]\ar@<-.8ex>[l]\ar@<.8ex>[l]
}
$$
Let
% original p. 413
$q\colon Z\to X$ and $r\colon T\to X$ be the canonical morphisms,
the notations $q', r', q'', r''$ having an obvious
meaning. By \Ref{XIII.1.11}~2), one has the following
essentially commutative diagram, whose rows are exact:
$$
\xymatrix{ f'^* i_*\Phi \ar[r]\ar[d]_a &
p'_*p'^*(f'^*i_*\Phi)\ar[d]_b\ar@<-.4ex>[r]\ar@<.4ex>[r] &
q'_*q'^*(f'^*i_*\Phi)\ar[d]_c\ar[r]\ar@<-.8ex>[r]\ar@<.8ex>[r]&
r'_*r'^*(f'^*i_*\Phi)\ar[d]_d \\ j_*f''^*\Phi \ar[r]&
j_*p''_*p''^*(f''^*\Phi)\ar@<-.4ex>[r]\ar@<.4ex>[r] &
j_*q''_*q''^*(f''^*\Phi)\ar[r]\ar@<-.8ex>[r]\ar@<.8ex>[r]&
j_*r''_*r''^*(f''^*\Phi) }
$$
Since $Y$ is smooth over~$S$, the morphism $fp$ is universally
locally $1$-aspherical for $\LL$ (SGA~XV~2.1),
and it follows from \cite[VII 2.1.7]{XIII.2} that $b$ is an equivalence of
categories. On the other hand, at the cost of restricting $S$ to a
nonempty open subset, one may assume that the morphisms $Z\to S$ and $T\to
S$ are universally locally $0$-acyclic. It follows
that the functors $c$ and $d$ are fully faithful, and the
diagram above then shows that $a$ is an equivalence, which
completes the proof.
\begin{corollary}
\label{XIII.3.4}
Let $k$ be a field of characteristic $p\geqslant0$, $p'$
the set of prime numbers distinct from~$p$, $f\colon X\to k$ a
coherent morphism. Suppose that one of the following two
conditions is satisfied:
\begin{enumerate}
\item[a)] $f$ is of finite type and the schemes of finite type of
dimension $\leqslant\dim X$ over an algebraic closure of~$k$
are desingularizable.
\item[b)] The schemes of finite type over an algebraic
closure of~$k$ are desingularizable.
\end{enumerate}
Then $f$ is universally locally $1$-aspherical for $p'$.
\end{corollary}

Case a) follows from~\Ref{XIII.3.3}. In case b), the question
being local on~$X$, one may assume $X$ affine; by passage
to the limit (SGA~4 XV~1.3), one reduces to the case where $X$ is
of finite type over~$k$.
\begin{corollary}
\label{XIII.3.5}
Let
% original p. 414
$S$ be an irreducible scheme with generic point
$s$, $f\colon X\to S$ a morphism of finite
presentation. Suppose that the schemes of finite type of dimension
$\leqslant\dim X_s$ over an algebraic closure $\overline{k}$
of~$\kres(s)$ are strongly desingularizable \textup{(SGA~5 I~3.1.5)}. If
$\LL$ denotes the set of prime numbers distinct from the
residue characteristics of~$S$, one can find a nonempty open
subset~$S_1$ of~$S$ such that, for every specialization
$\bar{s}_1\to\bar{s}_2$ of geometric points of~$S_1$, the
specialization morphism \eqref{XIII.2.10}
$$
\pi_1^{\LL}(X_{\bar{s}_1})\to\pi_1^{\LL}(X_{\bar{s}_2})
$$
is bijective.
\end{corollary}

By~\Ref{XIII.3.1} and~\Ref{XIII.3.3}, one may, at the cost of
restricting $S$ to a nonempty open subset, assume that $f$ is
locally $1$-aspherical for $\LL$, and that, for every finite constant sheaf
of~$\LL$-groups $F$ on~$X$, $(F,f)$ is
cohomologically proper in dimension $\leqslant 1$. It then follows
from~\Ref{XIII.1.14} that, for every specialization
$\bar{s}_1\to\bar{s}_2$ of geometric points of~$S_1$, the
specialization morphism
$$
\left(\R^1 f_* F\right)_{\bar{s}_2} \to \left(\R^1 f_*
F\right)_{\bar{s}_1}
$$
is bijective. The corollary is nothing but the translation of the
foregoing in terms of fundamental groups.
